% !TeX root = ../se200_referential_regimes.tex

% =============================================================================
\section{Reference Inventory}
\label{sec:inventory}
% =============================================================================

Before analyzing identity, we fix the carrier inventory.
An accountability substrate must let one state, inspect, and compare claims
about who is responsible, which authority governs, what happened,
where it applies, what was recorded, and what was acted upon.
It must do this without embedding causal or normative content at the substrate level.

The inventory is modest.
It is not a theorem, an upper ontology, or a complete domain ontology.
It states the carrier kinds required by the accountability setting analyzed here.
Those carriers are the referents that must remain addressable under the assumptions
of Section~\ref{sec:formal_setting}.

The inventory supplies identity carriers but does not determine identity regimes.
A reference kind identifies the referents that must remain addressable.
An identity regime specifies sameness for a carrier under transformation.
The lower-bound proof depends on that distinction.

% -----------------------------------------------------------------------------
\subsection{Reference Kinds and Identity Carriers}
\label{subsec:reference_kinds_identity_carriers}
% -----------------------------------------------------------------------------

\begin{definition}[Reference kind]
  \label{def:reference_kind}
  A \emph{reference kind} is a class of referents that a
  neutral accountability substrate must be able to identify, track, and
  distinguish from other referents without making object-level causal or normative commitments.
\end{definition}

\begin{definition}[Identity carrier]
  \label{def:identity_carrier}
  An \emph{identity carrier} is a referential position in the substrate to
  which an identity regime may be assigned.
  It is the object whose sameness, difference, and persistence are evaluated under transformation.
\end{definition}

Reference kinds supply identity carriers.
They do not themselves determine identity regimes.
A single reference kind may support more than one identity regime when
different identity bases respond differently to the same transformation.
Conversely, the fact that two carriers are both represented in the
substrate does not imply that they share an identity regime.
Regime assignment depends on transformation behavior, not on the name of the kind.

The inventory used in this paper is:

\begin{center}
  \begin{tabular}{@{}l>{\raggedright\arraybackslash}p{0.68\textwidth}@{}}
    \toprule
    Reference kind    & Required referent                                                             \\
    \midrule
    Obligation-bearer &
    a party that can bear responsibility across the occurrences
    discharging, violating, or contesting that responsibility                                         \\
    Rule              &
    a governing normative structure, distinct from the acts that enact, amend, apply, or challenge it \\
    Occurrence        &
    a happening individuated by temporal realization and provenance                                   \\
    Scope             &
    the applicability domain for a rule, including places, parties, periods,
    procedures, classifications, or conditions                                                        \\
    Record            &
    a descriptive indicator or record that carries reference without
    asserting causal or normative truth                                                               \\
    Plain referent    &
    a thing acted upon, regulated, located, transferred, or observed
    that neither bears obligations nor grounds authority as such                                      \\
    \bottomrule
  \end{tabular}
\end{center}

These six reference kinds are familiar from foundational ontology,
legal ontology, provenance, and institutional recordkeeping.
The paper does not claim novelty for the inventory itself.
The inventory matters because it supplies the carriers on which
the later identity analysis operates.

% -----------------------------------------------------------------------------
\subsection{Obligation-Bearers}
\label{subsec:inventory_obligation_bearers}
% -----------------------------------------------------------------------------

An accountability substrate must support reference to obligation-bearing entities.
An obligation-bearer is a party that can bear
responsibility, duty, liability, authority, or accountability
across occurrences in which those statuses are fulfilled, violated, transferred, contested, or discharged.
The substrate need not decide whether a party is responsible in a particular case.
It must preserve reference to the party that a framework may treat as responsible.

In multi-jurisdiction regulatory reporting, a corporate entity may be treated differently across jurisdictions.
One jurisdiction may individuate responsibility at the parent-entity level.
Another may individuate it at the subsidiary, facility, role, or delegated-agent level.
A neutral substrate need not settle which legal theory governs.
It must still allow the relevant bearer to remain referable across those competing frameworks.

Obligation-bearers cannot be reduced to the occurrences in which obligations are discharged.
The same bearer may participate in many occurrences.
The same occurrence may involve several bearers.
Collapsing bearer into occurrence destroys the ability to compare responsibility across time and across frameworks.

% -----------------------------------------------------------------------------
\subsection{Rules}
\label{subsec:inventory_rules}
% -----------------------------------------------------------------------------

An accountability substrate must support reference to governing rules or normative structures.
A rule is a structure that may be
enacted, amended, delegated, interpreted, applied, violated, or repealed.
The substrate need not assert that the rule is valid, justified, binding, or correctly interpreted.
It must preserve reference to the rule as the object of such attributed claims.

In legislative history tracking, a statute must remain referable
independently of the legislative act that enacted it, the amendment
that altered it, the agency action that applied it, and the judicial decision that interpreted it.
A statute can be contested while still remaining the same referent
for purposes of citation, comparison, and provenance.
If the rule is collapsed into its enactment or application occurrence,
then later disagreement about the occurrence changes reference to the rule itself.

Rules also preview one of the split cases below.
A rule may be individuated by imposed requirements
or by textual or internal structure.
That difference is not settled by naming the kind ``rule.''
It is settled by transformation behavior.

% -----------------------------------------------------------------------------
\subsection{Occurrences}
\label{subsec:inventory_occurrences}
% -----------------------------------------------------------------------------

An accountability substrate must support reference to \emph{occurrences}.
An occurrence is a happening individuated by temporal realization and provenance.
Illustrative examples include filings, inspections, votes, transfers, transactions,
observations, decisions, enactments, amendments, and enforcement actions.
The substrate must preserve reference to the occurrence as something that happened,
independently of how competing frameworks evaluate it.

In multi-jurisdiction regulatory reporting, two jurisdictions may disagree about
whether a filing constitutes compliance.
They must nevertheless be able to refer to the same filing.
In legislative history tracking, parties may disagree about the legal significance of an enactment.
They must nevertheless be able to refer to the same enactment occurrence.
In provenance analysis, institutions may disagree about the meaning or reliability of a measurement.
They must nevertheless be able to refer to the same measurement occurrence.

Occurrences cannot be reduced to obligation-bearers, rules, records, or scopes.
They are time-indexed referents.
They may instantiate, affect, or be described by other carriers, but they are not identical with them.

% -----------------------------------------------------------------------------
\subsection{Scopes}
\label{subsec:inventory_scopes}
% -----------------------------------------------------------------------------

An accountability substrate must support reference to scopes.
A scope is the domain where, to whom, or under what conditions a rule applies.
Scopes may be territorial, institutional, temporal, procedural, classificatory, jurisdictional, or conditional.
They may be nested, overlapping, partially coincident, or contested.

In multi-jurisdiction regulatory reporting, the same rule may apply to
different entities, places, facilities, transactions, or reporting periods under different frameworks.
If scope is represented as an attribute of the rule, then a change in
applicability can appear to change the identity of the rule.
If scope is represented as a property of obligation-bearers, then the
domain of applicability cannot be compared independently of the parties within it.
A neutral substrate needs scope as a separate referent.

Scopes also preview one of the split cases below.
A scope may be individuated extensionally, by covered cases,
or structurally, by internal jurisdictional or applicability organization.
Decomposition can preserve one basis while changing the other.
The reference kind ``scope'' therefore does not determine an identity regime.

% -----------------------------------------------------------------------------
\subsection{Records}
\label{subsec:inventory_records}
% -----------------------------------------------------------------------------

An accountability substrate must support reference to descriptive records.
A record is a descriptive indicator, document, dataset entry, observation record,
provenance item, or other referent that can be cited, transferred, compared, annotated, or analyzed without the substrate asserting a causal or normative claim.
The substrate may record that a source asserts a claim.
It does not thereby assert the claim.

In cross-institutional provenance, a record may be produced by one institution,
received by another, annotated by a third, and reanalyzed by a fourth.
The institutions may disagree about what the record shows or whether it supports a conclusion.
They must nevertheless be able to refer to the same record through custody, annotation, and reuse.

Records cannot be reduced to the entities, occurrences, rules, or claims they describe.
Stable reference to records is required independently of those referents,
including when a record supports an attributed causal or normative claim.

% -----------------------------------------------------------------------------
\subsection{Plain Referents}
\label{subsec:inventory_plain_referents}
% -----------------------------------------------------------------------------

An accountability substrate must also support reference to plain referents.
A plain referent is a thing acted upon, regulated, located, transferred, observed,
measured, or described that neither bears obligations nor grounds authority.
Examples include facilities, parcels, instruments, devices, environmental sites,
physical objects, infrastructure components, samples, and assets.
The category is intentionally modest.
It captures referents that are important to accountability without already being
obligation-bearers, rules, occurrences, scopes, or records.

Plain referents are needed because many accountability questions concern things
that are acted upon or regulated.
A regulated facility may be inspected.
A parcel may fall within a jurisdictional scope.
An instrument may produce a measurement record.
A sample may be transferred across institutions.
The substrate must preserve reference to the thing without deciding which
causal or normative interpretation applies to it.

Plain referents also preview one of the split cases below.
A referent may be individuated by locus,
such as the site or place that remains fixed.
It may also be individuated by the object occupying that locus.
Forking or replacement can separate those bases.
The reference kind ``plain referent'' therefore
does not by itself determine an identity regime.

% -----------------------------------------------------------------------------
\subsection{Summary}
\label{subsec:inventory_summary}
% -----------------------------------------------------------------------------

The inventory supplies six reference kinds:
obligation-bearer, rule, occurrence, scope, record, and plain referent.
These kinds are the carriers for the lower-bound construction
and do not determine identity regimes.

The inventory identifies the required referents.
The transformation analysis determines their behavior under ordinary record operations.
The core is derived after applying those operations
to the carrier-basis pairs.

The next section applies the transformation basis of
Section~\ref{sec:formal_setting} to these six carriers.
